\chapter{Entropion and Ectropion}

\section*{Definition}
Entropion is an inward rotation of the eyelid margin causing lashes to abrade the cornea and conjunctiva; ectropion is an outward sagging of the lower eyelid exposing the palpebral conjunctiva. Both are most commonly involutional (age-related) and affect the lower lid.

\section*{What Happens in Your Body}
Involutional entropion results from horizontal lid laxity, disinsertion of the lower lid retractors, and overriding of the preseptal orbicularis over the pretarsal orbicularis. Ectropion similarly stems from horizontal laxity combined with weakening of the lid retractors and canthal tendon attenuation, allowing the lid to evert.

\section*{Causes}
\begin{itemize}
  \item \textbf{Involutional (most common):} Age-related degeneration of lid support structures
  \item \textbf{Cicatricial:} Scarring of conjunctiva or posterior lamella from burns, surgery, or Stevens-Johnson syndrome
  \item \textbf{Spastic:} Sustained orbicularis contraction due to ocular surface irritation
  \item \textbf{Paralytic:} Facial nerve palsy (Bell's palsy, stroke) causing loss of orbicularis tone (ectropion)
  \item \textbf{Mechanical:} Tumor, edema, or prominent globe displacing the eyelid
\end{itemize}

\section*{When to See a Doctor}
See your doctor if:
\begin{itemize}
  \item Sensation of sand or gravel in the eye with chronic irritation
  \item Epiphora (excessive tearing) refractory to conservative measures
  \item Recurrent conjunctivitis or keratitis
  \item Visible inward or outward turning of the lower eyelid
  \item Corneal ulcer or opacity noted by you or caregiver
\end{itemize}

\section*{Related Symptoms}
\begin{itemize}
  \item Epiphora (tearing)
  \item Ocular irritation and foreign body sensation
  \item Conjunctival injection and erythematous eyelid margin
  \item Photophobia
  \item Mucoid discharge
\end{itemize}
\textbf{Important:} Untreated entropion with trichiasis can cause corneal scarring and vision loss, whereas chronic ectropion may lead to conjunctival keratinization and exposure keratopathy.

\section*{Self-Care / Home Management}
\begin{itemize}
  \item Artificial tears and lubricating ointment at bedtime for corneal protection
  \item Manual taping of the lower lid to temporarily correct mild ectropion
  \item Eyelid hygiene with warm compresses for associated blepharitis
  \item Temporary skin tape to evert the lid margin in spastic entropion
  \item Sunglasses to reduce exposure and photophobia
\end{itemize}

\section*{How Your Doctor Will Evaluate You}
\begin{itemize}
  \item External examination evaluating lid margin position relative to the globe in primary gaze
  \item Snap-back test and distraction test to quantify horizontal lid laxity
  \item Slit-lamp examination for corneal staining (fluorescein) and conjunctival changes
  \item Assessment of orbicularis strength and facial nerve function (if paralytic suspect)
  \item Eversion of the upper lid if entropion from conjunctival scarring suspected
\end{itemize}

\section*{Treatment Options}
\begin{itemize}
  \item Lubrication for mild cases without functional impairment
  \item Temporary botulinum toxin injection to weaken orbicularis overaction in spastic entropion
  \item Horizontal tightening procedures (lateral tarsal strip, lower lid retractor reinsertion) for involutional cases
  \item Everting sutures (Quickert or Wies procedure) for spastic or mild involutional entropion
  \item Posterior lamella grafting for cicatricial entropion
  \item Lid shortening with canthopexy or medial spindle procedure for ectropion repair
  \item Tarsorrhaphy for paralytic ectropion to protect the cornea
\end{itemize}

\section*{References}
\begin{thebibliography}{99}
\bibitem{entropion_ectropion:ref1} Damasceno RW, Osaki MH, Dantas PEC, et al. ``Involutional Entropion and Ectropion of the Lower Eyelid: A Review.'' Arq Bras Oftalmol. 2011;74(5):374--378.
\bibitem{entropion_ectropion:ref2} Bashour M, Harvey JT. ``Causes of Involutional Ectropion and Entropion: A Review.'' Ophthalmic Plast Reconstr Surg. 2000;16(1):47--53.
\bibitem{entropion_ectropion:ref3} Levine MR. ``Manual of Oculoplastic Surgery.'' 4th ed. Springer; 2018.
\bibitem{entropion_ectropion:ref4} Bartley GB. ``Acquired Ectropion and Entropion.'' In: Albert DM, Jakobiec FA, eds. Principles and Practice of Ophthalmology. 3rd ed. Saunders; 2008.
\bibitem{entropion_ectropion:ref5} Kanski JJ, Bowling B. ``Clinical Ophthalmology: A Systematic Approach.'' 7th ed. Elsevier; 2011.
\end{thebibliography}
